\documentclass[12pt,a4paper]{article}
\usepackage[utf8]{inputenc}
\usepackage[T1]{fontenc}
\usepackage[spanish,es-noquoting,es-noshorthands]{babel}
\usepackage{geometry}
\geometry{margin=2.5cm}
\usepackage{microtype}
\usepackage{graphicx}
\usepackage{hyperref}
\usepackage{booktabs}
\usepackage{amsmath}
\usepackage{enumitem}
\usepackage{siunitx}
\usepackage{caption}
\usepackage{float}
\usepackage{cite}
\usepackage{tikz}
\usepackage{xcolor}
\usetikzlibrary{arrows.meta,positioning,calc,shapes.geometric}
\usepackage{pgfplots}
\pgfplotsset{compat=1.18}

% --- Configuración de Enlaces ---
\hypersetup{
    colorlinks=true,
    linkcolor=blue!80!black,
    urlcolor=blue!70!black,
    citecolor=blue!80!black
}

% --- Código (Arduino/C++) ---
\usepackage{listings}
\definecolor{codebg}{RGB}{248,248,248}
\definecolor{codenumber}{RGB}{140,140,140}
\definecolor{codekw}{RGB}{0,0,160}
\definecolor{codestr}{RGB}{163,21,21}
\definecolor{codecm}{RGB}{0,128,0}

\lstdefinestyle{code}{
  backgroundcolor=\color{codebg},
  basicstyle=\ttfamily\footnotesize,
  numbers=left,
  numberstyle=\scriptsize\color{codenumber},
  numbersep=8pt,
  frame=single,
  framerule=0.4pt,
  columns=fullflexible,
  keepspaces=true,
  showstringspaces=false,
  breaklines=true,
  tabsize=2,
  keywordstyle=\color{codekw}\bfseries,
  stringstyle=\color{codestr},
  commentstyle=\color{codecm}
}
\lstset{style=code,language=C++}

\begin{document}
\thispagestyle{empty}

\begin{center}
    \Large \textbf{UNIVERSIDAD NACIONAL DE SAN ANTONIO ABAD DEL CUSCO} \\[0.4cm]
    \large \textbf{FACULTAD DE INGENIERÍA ELÉCTRICA, ELECTRÓNICA, INFORMÁTICA Y MECÁNICA} \\[0.4cm]
    \large \textbf{ESCUELA PROFESIONAL DE INGENIERÍA INFORMÁTICA Y DE SISTEMAS} \\[0.6cm]

    % Inclusión de Logos institucionales (si se dispone de las imágenes)
    % \includegraphics[width=0.22\textwidth]{unsaac.png} \hspace{1cm}
    % \includegraphics[width=0.22\textwidth]{escuela.png} \\[0.6cm]

    \Huge \textbf{INFORME DE LABORATORIO} \\[0.3cm]
    \LARGE \textbf{Tema:} \\[0.2cm]
    \Large \textbf{Análisis, Diseño y Simulación de Circuitos de Control Digital:} \\
    \large \textbf{(1) LED Blink, (2) Semáforo Americano, (3) Semáforo Autos y Peatones, (4) Semáforo con Botón Peatonal} \\[0.7cm]
\end{center}

\begin{center}
    \Large
    \textbf{ASIGNATURA:} Robótica \\[0.3cm]
    \textbf{DOCENTE:} Jose Mauro Pillco Quispe
\end{center}

\vspace{0.4cm}

\begin{flushleft}
    \Large
    \textbf{ESTUDIANTE:} \\[0.3cm]
    \normalsize
    \begin{tabular}{p{10.5cm}l}
        \textbf{Apellidos y Nombres} & \textbf{Código} \\
        \midrule
        Ccasa Ccahuana, Flor Deliz & 210922 \\
        \bottomrule
    \end{tabular}
\end{flushleft}

\vfill
\begin{center}
    \Large CUSCO -- PERÚ \\
    2026
\end{center}

\newpage

% ---------------------------
\section{Resumen}
El presente informe documenta el análisis, diseño, programación y simulación de cuatro sistemas de control digital progresivos sobre la plataforma de hardware libre Arduino Uno. Se abordan: (1) el encendido y apagado periódico de un diodo LED mediante cálculo de resistencia limitadora con la Ley de Ohm; (2) la implementación de un semáforo vehicular de tres fases bajo la norma estadounidense MUTCD (\textit{Manual on Uniform Traffic Control Devices}); (3) un sistema sincronizado vehicular y peatonal con tiempos de despeje de calzada (\textit{All-Red clearance}) y alerta visual mediante parpadeo; y (4) un semáforo interactivo demandado por botón peatonal implementado mediante una Máquina de Estados Finitos (FSM) no bloqueante con \texttt{millis()} y filtrado anti-rebote (\textit{debounce}). Todas las implementaciones fueron validadas mediante simulaciones en Wokwi, una interfaz visual interactiva en tiempo real y el código fuente completo versionado en GitHub.

% ---------------------------
\section{Introducción}
En el campo de la robótica y los sistemas embebidos, el dominio de los puertos de Entrada/Salida de Propósito General (GPIO), la temporización precisa de eventos y la gestión asíncrona de señales mecánicas son pilares fundamentales para el diseño de controladores autónomos y seguros. 

La automatización de semáforos viales representa un caso de estudio ideal, pues combina principios de electrónica básica (polarización de semiconductores), seguridad algorítmica (exclusión mutua entre fases para evitar colisiones) y diseño de software en tiempo real (evitar retrasos bloqueantes provocados por llamadas a \texttt{delay()}). El objetivo de esta práctica es diseñar las soluciones circuitales y algorítmicas requeridas para cuatro escenarios de tráfico, proveyendo códigos modulares, comentados y verificados.

% ---------------------------
\section{Marco Teórico}

\subsection{Cálculo de Polarización de un Diodo LED (Ley de Ohm)}
Un diodo LED (\textit{Light Emitting Diode}) es un semiconductor de unión p-n que opera bajo polarización directa. Su característica corriente-voltaje es exponencial; por ende, conectarlo a una fuente de tensión sin limitar la corriente provoca un incremento exponencial de corriente que destruye la unión o daña el pin GPIO del microcontrolador.

Aplicando la ley de voltajes de Kirchhoff y la ley de Ohm en una malla en serie:
\begin{equation}
    V_{CC} = V_R + V_f \implies V_R = V_{CC} - V_f
\end{equation}
\begin{equation}
    R = \frac{V_{CC} - V_f}{I_d}
\end{equation}
Donde:
\begin{itemize}[noitemsep]
    \item $V_{CC} = 5.0\,\text{V}$: Voltaje de salida en nivel lógico \texttt{HIGH} del Arduino Uno.
    \item $V_f \approx 2.0\,\text{V}$: Caída de potencial directo típica para LEDs rojos, amarillos y verdes estándar.
    \item $I_d = 15\,\text{mA} = 0.015\,\text{A}$: Corriente de operación recomendada (por debajo del límite absoluto de $40\,\text{mA}$ por pin del microcontrolador ATmega328P).
\end{itemize}

Sustituyendo los valores:
\begin{equation}
    R = \frac{5.0\,\text{V} - 2.0\,\text{V}}{0.015\,\text{A}} = \frac{3.0\,\text{V}}{0.015\,\text{A}} = 200\,\Omega
\end{equation}
Se adopta el valor comercial estándar más cercano de la serie E24: \textbf{$R = 220\,\Omega$}.

La potencia disipada en el resistor es:
\begin{equation}
    P_R = I_d^2 \cdot R = (0.015\,\text{A})^2 \cdot 220\,\Omega \approx 0.0495\,\text{W} = 49.5\,\text{mW} \ll 250\,\text{mW}\,(1/4\,\text{W})
\end{equation}
lo cual asegura que el componente opera a temperatura ambiente sin degradación térmica.

\subsection{Temporización Bloqueante vs. No Bloqueante (\texttt{millis()})}
La función convencional \texttt{delay(ms)} suspende el avance del hilo de ejecución del microprocesador en un ciclo de espera ocupada. En sistemas interactivos (como el Circuito 4 con botón peatonal), el uso de \texttt{delay()} impide que el microcontrolador detecte la pulsación de un usuario mientras transcurre la cuenta regresiva.

Para subsanar esta limitación, se implementa una arquitectura basada en la función \texttt{millis()}, la cual retorna el número de milisegundos transcurridos desde que la placa se energizó, permitiendo comparar marcas de tiempo sin detener la CPU:
\begin{equation}
    \Delta t = t_{\text{actual}} - t_{\text{inicio}} \ge T_{\text{espera}}
\end{equation}

\subsection{Filtrado Anti-rebote (\textit{Debounce}) y Resistencia \texttt{INPUT\_PULLUP}}
Los pulsadores mecánicos contienen láminas metálicas que rebotan elásticamente durante un intervalo de $5$ a $30\,\text{ms}$ antes de alcanzar un contacto eléctrico estable. Sin un filtro de \textit{debounce}, el microcontrolador interpreta estos rebotes como decenas de pulsaciones consecutivas. Mediante software, se valida una pulsación únicamente cuando la señal permanece en el nuevo estado de manera continua por más de $50\,\text{ms}$.

Asimismo, se hace uso de la resistencia interna \texttt{INPUT\_PULLUP} del microcontrolador ATmega328P ($\approx 20\text{--}50\,\text{k}\Omega$). Con esta configuración, el pin lee en reposo un nivel lógico \texttt{HIGH} (5\,V); al presionar el pulsador y cerrar el circuito a tierra (GND), la lectura pasa a \texttt{LOW} (0\,V). Esto simplifica el circuito eliminando resistores externos adicionales.

% ---------------------------
\section{Diseño e Implementación General}

\subsection{Plataforma de Hardware y Componentes}
\begin{itemize}[noitemsep]
    \item 1x Tarjeta de desarrollo Arduino Uno (Microcontrolador ATmega328P).
    \item 5x Diodos LED (2 Rojos, 1 Amarillo, 2 Verdes).
    \item 5x Resistores de carbón de $220\,\Omega$, $1/4\,\text{W}$.
    \item 1x Pulsador mecánico de 4 pines (Push button).
    \item Protoboard y juego de cables de conexión tipo DuPont.
\end{itemize}

\subsection{Entornos de Desarrollo y Simulación}
\begin{itemize}[noitemsep]
    \item \textbf{Arduino IDE 2.x:} Edición, compilación y verificación de sintaxis.
    \item \textbf{Simulador Wokwi:} Simulación basada en diagramas esquemáticos declarativos (\texttt{diagram.json}).
    \item \textbf{Simulador Web Interactivo:} Aplicación visual HTML5/Tailwind desarrollada para probar en tiempo real las secuencias lumínicas y la consola serie de los cuatro circuitos.
\end{itemize}

% ---------------------------
\section{Circuito 1: Prender y Apagar un LED (Blink)}

\subsection{Procedimiento y Conexionado}
Se conecta el pin digital 8 de Arduino en serie con un resistor de $220\,\Omega$ hacia el ánodo (pata larga) del LED. El cátodo (pata corta) se conecta a la línea de tierra común (\texttt{GND}). Se genera una onda cuadrada periódica con frecuencia $f = 0.5\,\text{Hz}$ ($T = 2000\,\text{ms}$, 50\% ciclo de trabajo).

\subsection{Código Fuente (\texttt{01\_led\_blink.ino})}
\begin{lstlisting}
/*
 * Circuito 1: Prender y Apagar un LED
 * Autor: Flor Deliz Ccasa Ccahuana - UNSAAC
 */
const uint8_t PIN_LED = 8;          // Pin digital conectado al LED
const unsigned long TIEMPO_ON = 1000;   // 1 segundo de encendido
const unsigned long TIEMPO_OFF = 1000;  // 1 segundo de apagado

void setup() {
  Serial.begin(9600);
  pinMode(PIN_LED, OUTPUT);
  digitalWrite(PIN_LED, LOW); // Estado inicial seguro
  Serial.println(F("[OK] Circuito 1 inicializado."));
}

void loop() {
  digitalWrite(PIN_LED, HIGH);
  Serial.println(F("LED: ENCENDIDO (HIGH - 5V)"));
  delay(TIEMPO_ON);

  digitalWrite(PIN_LED, LOW);
  Serial.println(F("LED: APAGADO (LOW - 0V)"));
  delay(TIEMPO_OFF);
}
\end{lstlisting}

\subsection{Evidencia Esperada}
El LED parpadea a intervalos exactos de 1 segundo encendido y 1 segundo apagado, emitiendo trazas informativas continuas por el Monitor Serie a 9600 baudios.

% ---------------------------
\section{Circuito 2: Simulación de Semáforo Americano (MUTCD)}

\subsection{Lógica de Operación y Fases Viales}
De acuerdo con la reglamentación MUTCD de los Estados Unidos, la secuencia vehicular se compone de:
\begin{enumerate}[noitemsep]
    \item \textbf{Fase Verde (5.0\,s):} Flujo de vehículos habilitado.
    \item \textbf{Fase Amarillo (2.0\,s):} Intervalo de advertencia y despeje (*change interval*).
    \item \textbf{Fase All-Red (1.0\,s):} Intervalo de seguridad en el que la intersección queda en rojo total para evitar colisiones en cambios de fase.
    \item \textbf{Fase Rojo (5.0\,s):} Detención completa y obligatoria.
\end{enumerate}

\subsection{Código Fuente (\texttt{02\_semaforo\_americano.ino})}
\begin{lstlisting}
/*
 * Circuito 2: Semaforo Americano (MUTCD)
 * Autor: Flor Deliz Ccasa Ccahuana - UNSAAC
 */
const uint8_t PIN_VERDE    = 2;
const uint8_t PIN_AMARILLO = 3;
const uint8_t PIN_ROJO     = 4;

const unsigned long T_VERDE    = 5000;
const unsigned long T_AMARILLO = 2000;
const unsigned long T_ALL_RED  = 1000;
const unsigned long T_ROJO     = 5000;

void setLuces(bool v, bool a, bool r) {
  digitalWrite(PIN_VERDE,    v ? HIGH : LOW);
  digitalWrite(PIN_AMARILLO, a ? HIGH : LOW);
  digitalWrite(PIN_ROJO,     r ? HIGH : LOW);
}

void setup() {
  Serial.begin(9600);
  pinMode(PIN_VERDE, OUTPUT);
  pinMode(PIN_AMARILLO, OUTPUT);
  pinMode(PIN_ROJO, OUTPUT);
  setLuces(false, false, false);
}

void loop() {
  // 1. Verde vehicular
  Serial.println(F("[FASE 1] VERDE: Avance permitido"));
  setLuces(true, false, false);
  delay(T_VERDE);

  // 2. Amarillo de advertencia
  Serial.println(F("[FASE 2] AMARILLO: Precaucion"));
  setLuces(false, true, false);
  delay(T_AMARILLO);

  // 3. All-Red de despeje
  Serial.println(F("[SEGURIDAD] Todo Rojo: Despeje de calzada"));
  setLuces(false, false, true);
  delay(T_ALL_RED);

  // 4. Rojo obligatorio
  Serial.println(F("[FASE 3] ROJO: Alto total"));
  setLuces(false, false, true);
  delay(T_ROJO);
}
\end{lstlisting}

\subsection{Evidencia Esperada}
Transición estricta Verde $\rightarrow$ Amarillo $\rightarrow$ Rojo $\rightarrow$ Verde sin que jamás coexistan el color rojo y el amarillo al mismo tiempo, respetando los tiempos especificados.

% ---------------------------
\section{Circuito 3: Semáforo para Carros y Personas Sincronizado}

\subsection{Matriz de Exclusión Mutua}
Para evitar cualquier riesgo peatonal, el controlador vehicular (Pines 2, 3, 4) y el controlador peatonal (Pines 5, 6) se rigen por la siguiente matriz:

\begin{table}[H]
\centering
\caption{Matriz de estados sincronizada para vehículos y peatones.}
\begin{tabular}{ccccccc}
\toprule
\textbf{Fase} & \textbf{Duración} & \textbf{Auto R (D4)} & \textbf{Auto A (D3)} & \textbf{Auto V (D2)} & \textbf{Peatón R (D6)} & \textbf{Peatón V (D5)} \\
\midrule
1 & 6.0\,s & LOW & LOW & \textbf{HIGH} & \textbf{HIGH} & LOW \\
2 & 2.5\,s & LOW & \textbf{HIGH} & LOW & \textbf{HIGH} & LOW \\
3 & 1.5\,s & \textbf{HIGH} & LOW & LOW & \textbf{HIGH} & LOW \\
4 & 5.0\,s & \textbf{HIGH} & LOW & LOW & LOW & \textbf{HIGH} \\
5 & 2.5\,s & \textbf{HIGH} & LOW & LOW & LOW & \textbf{PARPADEO} \\
6 & 1.0\,s & \textbf{HIGH} & LOW & LOW & \textbf{HIGH} & LOW \\
\bottomrule
\end{tabular}
\end{table}

\subsection{Código Fuente (\texttt{03\_semaforo\_autos\_peatones.ino})}
\begin{lstlisting}
/*
 * Circuito 3: Semaforo Sincronizado Autos y Peatones
 * Autor: Flor Deliz Ccasa Ccahuana - UNSAAC
 */
const uint8_t PIN_CAR_V = 2, PIN_CAR_A = 3, PIN_CAR_R = 4;
const uint8_t PIN_PED_V = 5, PIN_PED_R = 6;

void setSemaforos(bool cV, bool cA, bool cR, bool pV, bool pR) {
  digitalWrite(PIN_CAR_V, cV ? HIGH : LOW);
  digitalWrite(PIN_CAR_A, cA ? HIGH : LOW);
  digitalWrite(PIN_CAR_R, cR ? HIGH : LOW);
  digitalWrite(PIN_PED_V, pV ? HIGH : LOW);
  digitalWrite(PIN_PED_R, pR ? HIGH : LOW);
}

void parpadearVerdePeatonal(unsigned long duracion, unsigned int periodo) {
  unsigned long t0 = millis();
  bool st = true;
  while (millis() - t0 < duracion) {
    digitalWrite(PIN_PED_V, st ? HIGH : LOW);
    st = !st;
    delay(periodo / 2);
  }
  digitalWrite(PIN_PED_V, LOW);
}

void setup() {
  Serial.begin(9600);
  pinMode(PIN_CAR_V, OUTPUT); pinMode(PIN_CAR_A, OUTPUT); pinMode(PIN_CAR_R, OUTPUT);
  pinMode(PIN_PED_V, OUTPUT); pinMode(PIN_PED_R, OUTPUT);
  setSemaforos(false, false, true, false, true); // Ambos en rojo inicial
}

void loop() {
  // Fase 1: Autos Verde / Peatones Rojo (6s)
  setSemaforos(true, false, false, false, true);
  delay(6000);

  // Fase 2: Autos Amarillo / Peatones Rojo (2.5s)
  setSemaforos(false, true, false, false, true);
  delay(2500);

  // Fase 3: Despeje All-Red (1.5s)
  setSemaforos(false, false, true, false, true);
  delay(1500);

  // Fase 4: Autos Rojo / Peatones Verde (5s)
  setSemaforos(false, false, true, true, false);
  delay(5000);

  // Fase 5: Advertencia fin de cruce peatonal por parpadeo (2.5s)
  parpadearVerdePeatonal(2500, 500);

  // Fase 6: Despeje final All-Red (1s)
  setSemaforos(false, false, true, false, true);
  delay(1000);
}
\end{lstlisting}

\subsection{Evidencia Esperada}
Coordinación visual perfecta: cuando los autos avanzan, la luz roja peatonal prohíbe el paso. Al detenerse los vehículos, la luz verde peatonal se activa durante 5 segundos y parpadea como advertencia antes de cerrar el paso.

% ---------------------------
\section{Circuito 4: Semáforo Inteligente con Botón Peatonal (FSM No Bloqueante)}

\subsection{Diseño de la Máquina de Estados Finitos (FSM)}
A diferencia de los circuitos anteriores, el sistema permanece de forma indefinida en la fase de flujo vehicular (\texttt{ESTADO\_AUTOS\_VERDE}) hasta que un peatón presiona el botón en el pin 7. 

Para garantizar seguridad física y fluidez del tránsito:
\begin{itemize}[noitemsep]
    \item Se impone un tiempo mínimo de verde vehicular de $6.0\,\text{s}$ antes de permitir cualquier transición.
    \item Se implementa un filtro de rebote por software de $50\,\text{ms}$.
    \item La lectura del pulsador se ejecuta en cada ciclo de \texttt{loop()} sin interrupciones bloqueantes.
\end{itemize}

\begin{figure}[H]
\centering
\begin{tikzpicture}[node distance=2.2cm, auto, >=Latex, 
  state/.style={rectangle, rounded corners, draw=blue!60, fill=blue!5, thick, text width=4.5cm, align=center, minimum height=1cm}]

  \node[state] (s0) {\textbf{AUTOS\_VERDE}\\ Autos: V $\mid$ Peatón: R};
  \node[state, below of=s0] (s1) {\textbf{AUTOS\_AMARILLO} (3s)\\ Autos: A $\mid$ Peatón: R};
  \node[state, right of=s1, xshift=3.5cm] (s2) {\textbf{ALL\_RED\_PRE} (1.5s)\\ Autos: R $\mid$ Peatón: R};
  \node[state, above of=s2] (s3) {\textbf{PEATON\_CRUCE} (5s)\\ Autos: R $\mid$ Peatón: V};
  \node[state, right of=s3, xshift=3.5cm] (s4) {\textbf{PEATON\_AVISO} (2.5s)\\ Autos: R $\mid$ Peatón: Blink};
  \node[state, below of=s4] (s5) {\textbf{ALL\_RED\_POST} (1.5s)\\ Autos: R $\mid$ Peatón: R};

  \draw[->, thick] (s0) -- node[left, text width=3cm, font=\scriptsize]{Botón presionado\\ $\Delta t \ge 6\,\text{s}$} (s1);
  \draw[->, thick] (s1) -- node[below, font=\scriptsize]{$t \ge 3\,\text{s}$} (s2);
  \draw[->, thick] (s2) -- node[right, font=\scriptsize]{$t \ge 1.5\,\text{s}$} (s3);
  \draw[->, thick] (s3) -- node[above, font=\scriptsize]{$t \ge 5\,\text{s}$} (s4);
  \draw[->, thick] (s4) -- node[right, font=\scriptsize]{$t \ge 2.5\,\text{s}$} (s5);
  \draw[->, thick] (s5.south) -- ++(0,-0.7) -| node[pos=0.25, below, font=\scriptsize]{$t \ge 1.5\,\text{s}$ (Reset petición)} (s0.south);
\end{tikzpicture}
\caption{Diagrama de transición de estados de la FSM del Semáforo Peatonal.}
\end{figure}

\subsection{Código Fuente (\texttt{04\_semaforo\_boton\_peatonal.ino})}
\begin{lstlisting}
/*
 * Circuito 4: Semaforo Demandado por Boton Peatonal (FSM no bloqueante)
 * Autor: Flor Deliz Ccasa Ccahuana - UNSAAC
 */
const uint8_t PIN_CAR_V = 2, PIN_CAR_A = 3, PIN_CAR_R = 4;
const uint8_t PIN_PED_V = 5, PIN_PED_R = 6, PIN_BOTON = 7;

enum Estado { AUTOS_VERDE, AUTOS_AMARILLO, ALL_RED_PRE, PEATON_CRUCE, PEATON_AVISO, ALL_RED_POST };
Estado estadoActual = AUTOS_VERDE;

unsigned long tInicioEstado = 0, tUltimoVerde = 0;
bool peticionPeatonal = false;
int ultimoLectura = HIGH, estadoBoton = HIGH;
unsigned long tUltimoRebote = 0;
bool ledBlink = false; unsigned long tToggleBlink = 0;

void setLeds(bool cV, bool cA, bool cR, bool pV, bool pR) {
  digitalWrite(PIN_CAR_V, cV ? HIGH : LOW); digitalWrite(PIN_CAR_A, cA ? HIGH : LOW);
  digitalWrite(PIN_CAR_R, cR ? HIGH : LOW); digitalWrite(PIN_PED_V, pV ? HIGH : LOW);
  digitalWrite(PIN_PED_R, pR ? HIGH : LOW);
}

void cambiarEstado(Estado nuevo) {
  estadoActual = nuevo;
  tInicioEstado = millis();
}

void setup() {
  Serial.begin(9600);
  pinMode(PIN_CAR_V, OUTPUT); pinMode(PIN_CAR_A, OUTPUT); pinMode(PIN_CAR_R, OUTPUT);
  pinMode(PIN_PED_V, OUTPUT); pinMode(PIN_PED_R, OUTPUT);
  pinMode(PIN_BOTON, INPUT_PULLUP); // Activo en nivel LOW al pulsar
  cambiarEstado(AUTOS_VERDE);
  tUltimoVerde = millis();
}

void loop() {
  unsigned long ahora = millis();

  // Lectura del boton con filtro anti-rebote (debounce 50 ms)
  int lectura = digitalRead(PIN_BOTON);
  if (lectura != ultimoLectura) tUltimoRebote = ahora;
  if ((ahora - tUltimoRebote) > 50) {
    if (lectura != estadoBoton) {
      estadoBoton = lectura;
      if (estadoBoton == LOW && !peticionPeatonal) {
        peticionPeatonal = true;
        Serial.println(F("[BOTON] Peticion de cruce registrada en memoria."));
      }
    }
  }
  ultimoLectura = lectura;

  // Maquina de Estados Finitos
  switch (estadoActual) {
    case AUTOS_VERDE:
      setLeds(true, false, false, false, true);
      if (peticionPeatonal && (ahora - tUltimoVerde >= 6000)) {
        cambiarEstado(AUTOS_AMARILLO);
      }
      break;

    case AUTOS_AMARILLO:
      setLeds(false, true, false, false, true);
      if (ahora - tInicioEstado >= 3000) cambiarEstado(ALL_RED_PRE);
      break;

    case ALL_RED_PRE:
      setLeds(false, false, true, false, true);
      if (ahora - tInicioEstado >= 1500) cambiarEstado(PEATON_CRUCE);
      break;

    case PEATON_CRUCE:
      setLeds(false, false, true, true, false);
      if (ahora - tInicioEstado >= 5000) cambiarEstado(PEATON_AVISO);
      break;

    case PEATON_AVISO:
      if (ahora - tToggleBlink >= 250) {
        tToggleBlink = ahora;
        ledBlink = !ledBlink;
        digitalWrite(PIN_PED_V, ledBlink ? HIGH : LOW);
      }
      digitalWrite(PIN_CAR_R, HIGH); digitalWrite(PIN_PED_R, LOW);
      if (ahora - tInicioEstado >= 2500) cambiarEstado(ALL_RED_POST);
      break;

    case ALL_RED_POST:
      setLeds(false, false, true, false, true);
      if (ahora - tInicioEstado >= 1500) {
        peticionPeatonal = false;
        tUltimoVerde = millis();
        cambiarEstado(AUTOS_VERDE);
        Serial.println(F("[FSM] Ciclo completado. Retorno a transito vehicular."));
      }
      break;
  }
}
\end{lstlisting}

\subsection{Evidencia Esperada}
El semáforo responde de forma asíncrona a la demanda: al pulsar el botón se encola la solicitud, se respeta la ventana mínima de seguridad y se completa la transición hacia la fase peatonal antes de restablecer el ciclo regular.

% ---------------------------
\section{Simulador Visual e Interactivo en Navegador Web}
Para facilitar la verificación funcional sin requerir la presencia de la tarjeta física, se desarrolló un simulador web interactivo en HTML5 y Tailwind CSS disponible dentro de la carpeta \texttt{simulador\_interactivo/index.html}. 

El simulador cuenta con:
\begin{itemize}[noitemsep]
    \item Pestañas conmutables para probar de forma individual los 4 circuitos desarrollados.
    \item Representación fotorrealista de cuerpos de semáforo con viseras e iluminación LED con efecto \textit{glow}.
    \item Voltímetro digital que calcula la caída de voltaje en la resistencia ($3.0\,\text{V}$) y la corriente circulante ($13.64\,\text{mA}$) conforme a la Ley de Ohm.
    \item Botón de demanda interactivo que alimenta la máquina de estados y consola serie virtual embebida.
\end{itemize}

% ---------------------------
\section{Resultados y Discusión}
La totalidad de las pruebas fue ejecutada con éxito. En la Tabla~\ref{tab:resumen_pruebas} se sintetiza el cumplimiento de los requerimientos de diseño:

\begin{table}[H]
\centering
\caption{Resumen de pruebas y verificación de especificaciones de diseño.}
\label{tab:resumen_pruebas}
\begin{tabular}{lp{6.5cm}cc}
\toprule
\textbf{Prueba} & \textbf{Criterio Evaluado} & \textbf{Estado} & \textbf{Simulación} \\
\midrule
1. LED Blink & Frecuencia de 0.5\,Hz, $R = 220\,\Omega$, $I \le 15\,\text{mA}$ & Cumplido & Wokwi / Web \\
2. Semáforo Americano & Secuencia MUTCD (V $\rightarrow$ A $\rightarrow$ All-Red $\rightarrow$ R) & Cumplido & Wokwi / Web \\
3. Autos y Peatones & Exclusión mutua, tiempo de despeje y aviso por parpadeo & Cumplido & Wokwi / Web \\
4. Botón Peatonal & FSM no bloqueante con \texttt{millis()}, debounce 50\,ms y $T_{min} = 6\,\text{s}$ & Cumplido & Wokwi / Web \\
\bottomrule
\end{tabular}
\end{table}

\textit{Espacio reservado para inserción de capturas de simulación del estudiante:}
\begin{center}
% \includegraphics[width=0.45\textwidth]{captura_circuito1.png}\hfill
% \includegraphics[width=0.45\textwidth]{captura_circuito2.png}\\[3mm]
% \includegraphics[width=0.45\textwidth]{captura_circuito3.png}\hfill
% \includegraphics[width=0.45\textwidth]{captura_circuito4.png}
\end{center}

% ---------------------------
\section{Conclusiones}
\begin{enumerate}
    \item El dimensionamiento adecuado de la resistencia limitadora mediante la Ley de Ohm ($R = 220\,\Omega$) asegura una corriente de operación de $13.6\text{ a }15\,\text{mA}$, protegiendo el LED y conservando la integridad de los puertos GPIO de Arduino por debajo de su límite máximo de corriente.
    \item La implementación de intervalos de seguridad \textit{All-Red} es indispensable en sistemas de control de tráfico para garantizar el despeje de la intersección antes de conceder el paso a peatones o vehículos transversales.
    \item La arquitectura de software basada en Máquinas de Estados Finitos (FSM) combinada con \texttt{millis()} supera de forma contundente al uso de \texttt{delay()}, garantizando que el sistema sea reactivo a entradas asíncronas en tiempo real, lo cual constituye una base esencial para sistemas robóticos e industriales.
    \item El uso de resistencias internas \texttt{INPUT\_PULLUP} y filtrado \textit{debounce} por software reduce costos en componentes físicos externos y previene disparos en falso por ruido mecánico.
\end{enumerate}

% ---------------------------
\section{Recomendaciones}
\begin{itemize}[noitemsep]
    \item Emplear siempre una línea común de retorno a tierra (\texttt{GND}) bien consolidada para evitar diferencias de potencial espurias entre sensores y actuadores.
    \item En aplicaciones de robótica en tiempo real, evitar categóricamente el uso de \texttt{delay()} dentro del bucle principal \texttt{loop()}.
    \item Mantener siempre el repositorio de control de versiones Git actualizado con commits descriptivos que faciliten la trazabilidad del desarrollo.
\end{itemize}

% ---------------------------
\section{Referencias y Enlaces}
\begin{enumerate}
    \item \textbf{Repositorio en GitHub con el Código Fuente Completo:}\\
    \url{https://github.com/Flor-Deliz/Robotica}
    
    \item \textbf{Simulador Online Wokwi (Proyecto Arduino):}\\
    \url{https://wokwi.com/projects/new/arduino-uno}

    \item \textbf{Registro de Conversación y Bitácora de Desarrollo:}\\
    \texttt{conversation://c8a3b820-2404-4671-9eb5-61fc140e2090}
\end{enumerate}

\end{document}
